\section{The truncated expansion as an inclusion--exclusion sum}
\label{sec:hastings_cluster_truncation}

% Hastings, arXiv:cond-mat/0508554v3, proj3new.tex, lines 1090--1138 (active
% bonds, clusters, the boundary) and 1197--1225 (the inclusion--exclusion sum
% and the remark added).

Expanding $e^{-\beta H}=\sum_n\frac{(-\beta)^n}{n!}\bigl(\sum_kH_k\bigr)^n$ into
words $H_{k_1}\cdots H_{k_n}$, the truncated operator
$\rho(\beta,l_{\mathrm{proj}})$ keeps words containing no cluster of active
bonds of diameter larger than $l_{\mathrm{proj}}$
\cite[lines 1090--1114]{Hastings2006Solving} and equals the inclusion--exclusion
sum $\rho-\rho_1+\rho_2-\cdots$ over families of large clusters
\cite[lines 1218--1224]{Hastings2006Solving}, provided the cluster boundary
consists of terms whose supports meet the cluster support
\cite[lines 1131--1134]{Hastings2006Solving},
\cite{gap:hastings05_cluster_boundary}.

\begin{definition}[Clusters of a set of terms]
    \label{def:hastings_cluster_of_terms}
    \lean{TNLean.Lattice.ClusterExpansion.overlapGraph,
        TNLean.Lattice.ClusterExpansion.boundary,
        TNLean.Lattice.ClusterExpansion.mem_boundary_overlapGraph,
        TNLean.Lattice.ClusterExpansion.IsCluster,
        TNLean.Lattice.ClusterExpansion.IsAdmissible}
    \leanok
    \uses{thm:hastings_cluster_boundary_count}
    Let $G$ be a graph on a finite term index set $K$ (in the overlap graph,
    $k\sim k'$ when $k\ne k'$ and their supports meet). For $C\subseteq K$, its
    \emph{boundary} is $\partial C=\{k\in K\setminus C:\exists c\in C,\ k\sim c\}$.
    A \emph{cluster} of $A\subseteq K$ is a $G$-connected subset $C\subseteq A$
    with $A\cap\partial C=\emptyset$. A family $U$ of subsets of $K$ is
    \emph{admissible} if $C\cap\partial C'=\emptyset$ for all $C,C'\in U$.
\end{definition}

\begin{lemma}[Clusters that meet]
    \label{thm:hastings_clusters_meet}
    \lean{TNLean.Lattice.ClusterExpansion.subset_of_connected_of_disjoint_boundary,
        TNLean.Lattice.ClusterExpansion.eq_of_isCluster_of_inter_nonempty,
        TNLean.Lattice.ClusterExpansion.pairwiseDisjoint_of_admissible,
        TNLean.Lattice.ClusterExpansion.biUnion_boundary_eq_of_admissible,
        TNLean.Lattice.ClusterExpansion.card_biUnion_of_admissible,
        TNLean.Lattice.ClusterExpansion.card_add_card_biUnion_boundary_le}
    \leanok
    \uses{def:hastings_cluster_of_terms}
    Intersecting clusters of $A$ coincide. For an admissible family $U$ of
    connected sets with $C_U=\bigcup_{C\in U}C$, the members of $U$ are pairwise
    disjoint, $\partial C_U=\bigcup_{C\in U}\partial C$, and if
    $n_C+n_{\partial C}\le Sn_C$ for all $C\in U$, then
    $n_{C_U}+n_{\partial C_U}\le Sn_{C_U}$.
\end{lemma}
\begin{proof}
    \leanok
    If $C$ is connected, $C\cap C'\ne\emptyset$, and $C\cap\partial C'=\emptyset$,
    any path in $C$ from $C\cap C'$ stays in $C'$, so $C\subseteq C'$; symmetry
    gives equality for intersecting clusters of $A$ and pairwise disjointness for
    admissible families, whence $\partial C_U=\bigcup_{C\in U}\partial C$ and
    the additive size bound.
\end{proof}

\begin{lemma}[Words containing a set of terms]
    \label{thm:hastings_words_inclusion_exclusion}
    \lean{TNLean.Lattice.ClusterExpansion.sum_pow_eq_sum_words,
        TNLean.Lattice.ClusterExpansion.sum_powerset_neg_one_pow_card_disjoint,
        TNLean.Lattice.ClusterExpansion.sum_words_contain_avoid}
    \leanok
    In any ring, $(\sum_{k\in s}a_k)^n=\sum_{w\in s^n}a_{w_1}\cdots a_{w_n}$,
    and for disjoint $B,C\subseteq K$, the sum of $a_{w_1}\cdots a_{w_n}$ over
    words of length $n$ whose letter set $X$ satisfies $C\subseteq X$ and
    $X\cap B=\emptyset$ equals
    \begin{align}
        \sum_{D\subseteq C}(-1)^{|D|}\Bigl(\sum_{k\notin B\cup D}a_k\Bigr)^n.
        \label{eq:hastings_words_inclusion_exclusion}
    \end{align}
\end{lemma}
\begin{proof}
    \leanok
    For $X\cap B=\emptyset$,
    $\sum_{D\subseteq C\setminus X}(-1)^{|D|}=[C\subseteq X]$.
\end{proof}

\begin{lemma}[The cluster part as a sum of words]
    \label{thm:hastings_cluster_part_words}
    \lean{TNLean.Lattice.ClusterExpansion.hasSum_div_factorial_smul_pow,
        TNLean.Lattice.ClusterExpansion.hasSum_words_clusterPart}
    \leanok
    \uses{thm:hastings_words_inclusion_exclusion, def:hastings_cluster_part}
    For $H=\sum_{k\in K}H_k$ and disjoint $B,C\subseteq K$, the words whose
    letter set contains $C$ and avoids $B$, weighted by $(-\beta)^n/n!$, sum to
    $\rho(\beta,C)$ of (\ref{eq:hastings_cluster_part}) with boundary $B$.
\end{lemma}
\begin{proof}
    \leanok
    Sum (\ref{eq:hastings_words_inclusion_exclusion}) weighted by
    $(-\beta)^n/n!$ over $n\ge0$.
\end{proof}

\begin{definition}[The truncated operator]
    \label{def:hastings_truncated_expansion}
    \lean{TNLean.Lattice.ClusterExpansion.NoLargeCluster,
        TNLean.Lattice.ClusterExpansion.largeConnected,
        TNLean.Lattice.ClusterExpansion.truncatedExpTerm,
        TNLean.Lattice.ClusterExpansion.truncatedExp}
    \leanok
    \uses{def:hastings_cluster_of_terms}
    For a set $T$ of large connected subsets of $K$, $\rho(\beta,l)$ is the sum
    over $n\ge0$ of $\frac{(-\beta)^n}{n!}\sum H_{k_1}\cdots H_{k_n}$ restricted
    to words whose letter set has no cluster in $T$
    \cite[lines 1101--1104]{Hastings2006Solving}.
\end{definition}

\begin{lemma}[Inclusion and exclusion for one word]
    \label{thm:hastings_word_large_clusters}
    \lean{TNLean.Lattice.ClusterExpansion.mem_largeConnected,
        TNLean.Lattice.ClusterExpansion.forall_isCluster_iff,
        TNLean.Lattice.ClusterExpansion.ite_noLargeCluster_eq_sum}
    \leanok
    \uses{thm:hastings_clusters_meet, def:hastings_truncated_expansion}
    For $A\subseteq K$, summing over admissible $U\subseteq T$ with
    $\bigcup_{C\in U}C\subseteq A$ and
    $A\cap\bigcup_{C\in U}\partial C=\emptyset$ gives
    \begin{align}
        [\text{no cluster of }A\text{ lies in }T]
        =\sum_{U}(-1)^{|U|}.
        \label{eq:hastings_word_large_clusters}
    \end{align}
\end{lemma}
\begin{proof}
    \leanok
    Such $U$ are precisely the subsets of the set $L$ of large clusters of $A$,
    and $\sum_{U\subseteq L}(-1)^{|U|}=[L=\emptyset]$.
\end{proof}

\begin{lemma}[The truncated operator by inclusion and exclusion]
    \label{thm:hastings_truncation_inclusion_exclusion}
    \lean{TNLean.Lattice.ClusterExpansion.hasSum_truncatedExpTerm,
        TNLean.Lattice.ClusterExpansion.truncatedExp_eq_sum,
        TNLean.Lattice.ClusterExpansion.exp_sub_truncatedExp}
    \leanok
    \uses{thm:hastings_word_large_clusters, thm:hastings_cluster_part_words,
        def:hastings_truncated_expansion}
    For $H=\sum_{k\in K}H_k$, the series $\rho(\beta,l)$ converges and, summing
    over admissible $U\subseteq T$ with $\rho(\beta,U)$ given by
    (\ref{eq:hastings_cluster_part}) for $\bigl(\bigcup_{C\in U}C,\bigcup_{C\in U}\partial C\bigr)$,
    \begin{align}
        \rho(\beta,l)=\sum_U(-1)^{|U|}\rho(\beta,U),\qquad
        e^{-\beta H}-\rho(\beta,l)=-\sum_{U\neq\emptyset}(-1)^{|U|}\rho(\beta,U)
        \label{eq:hastings_truncation_inclusion_exclusion}
    \end{align}
    \cite[lines 1197--1225]{Hastings2006Solving},
    \cite{gap:hastings05_cluster_boundary}.
\end{lemma}
\begin{proof}
    \leanok
    Substitute (\ref{eq:hastings_word_large_clusters}) into
    Definition~\ref{def:hastings_truncated_expansion}, exchange sums, and apply
    Lemma~\ref{thm:hastings_cluster_part_words} (with $\rho(\beta,\emptyset)=e^{-\beta H}$).
\end{proof}

\begin{lemma}[The bound summed by the percolation estimate]
    \label{thm:hastings_truncation_trace_bound}
    \lean{TNLean.Lattice.ClusterExpansion.norm_trace_mul_exp_sub_truncatedExp_le}
    \leanok
    \uses{thm:hastings_truncation_inclusion_exclusion,
        thm:hastings_cluster_part_several}
    If $H=\sum_kH_k$ has Hermitian terms with $\|H_k\|\le J$, $\beta\ge0$, and
    $n_C+n_{\partial C}\le Sn_C$ for all $C\in T$, then for every $O$, with
    $y=(e^{J\beta}-1)e^{SJ\beta}$ and $Z=\operatorname{tr}e^{-\beta H}$,
    \begin{align}
        \bigl|\operatorname{tr}\bigl(O(e^{-\beta H}-\rho(\beta,l))\bigr)\bigr|
        \le\|O\|Z\sum_{U\neq\emptyset}\prod_{C\in U}y^{n_C}.
        \label{eq:hastings_truncation_trace_bound}
    \end{align}
\end{lemma}
\begin{proof}
    \leanok
    \uses{thm:hastings_truncation_inclusion_exclusion,
        thm:hastings_cluster_part_several, thm:hastings_clusters_meet}
    Apply Lemmas~\ref{thm:hastings_cluster_part_several}
    and~\ref{thm:hastings_clusters_meet} to each term of
    (\ref{eq:hastings_truncation_inclusion_exclusion}).
\end{proof}

\begin{definition}[The operator of a set of terms]
    \label{def:hastings_cluster_operator}
    \lean{TNLean.Lattice.ClusterExpansion.clusterOp,
        TNLean.Lattice.ClusterExpansion.clusterOp_eq_clusterPart,
        TNLean.Lattice.ClusterExpansion.hasSum_words_clusterOp}
    \leanok
    \uses{thm:hastings_cluster_part_words}
    For $A\subseteq K$ and $H_D=\sum_{k\in D}H_k$, define
    \begin{align}
        E_A=\sum_{D\subseteq A}(-1)^{|A\setminus D|}e^{-\beta H_D},
        \label{eq:hastings_cluster_operator}
    \end{align}
    which sums the words in the expansion of $e^{-\beta H}$ whose letter set is
    exactly $A$ \cite[lines 1268--1269]{Hastings2006Solving}.
\end{definition}

\begin{lemma}[The truncated operator as a sum of products over clusters]
    \label{thm:hastings_truncation_cluster_products}
    \lean{TNLean.Lattice.ClusterExpansion.truncatedExp_eq_sum_clusterOp,
        TNLean.Lattice.ClusterExpansion.clusterOp_union,
        TNLean.Lattice.ClusterExpansion.clusters,
        TNLean.Lattice.ClusterExpansion.biUnion_clusters,
        TNLean.Lattice.ClusterExpansion.clusterOp_eq_noncommProd_clusters,
        TNLean.Lattice.ClusterExpansion.truncatedExp_eq_sum_noncommProd_clusters,
        TNLean.Lattice.ClusterExpansion.commute_of_not_adj_overlapGraph,
        TNLean.Lattice.ClusterExpansion.truncatedExp_eq_sum_noncommProd_clusters_overlapGraph}
    \leanok
    \uses{def:hastings_cluster_operator, def:hastings_truncated_expansion,
        def:hastings_cluster_of_terms}
    Summing over $A\subseteq K$ with no cluster in $T$,
    $\rho(\beta,l)=\sum_AE_A$. If non-adjacent terms commute (in particular for
    the overlap graph when each $H_k$ acts on $S_k$), then
    $E_A=\prod_{C\in\mathrm{clusters}(A)}E_C$ with commuting factors
    \cite[lines 1108--1110]{Hastings2006Solving}.
\end{lemma}
\begin{proof}
    \leanok
    \uses{def:hastings_cluster_operator, thm:hastings_clusters_meet}
    Grouping words by letter set $A$ gives $\rho(\beta,l)=\sum_AE_A$. For
    disjoint commuting sets $A,B$,
    $e^{-\beta H_{D_1\cup D_2}}=e^{-\beta H_{D_1}}e^{-\beta H_{D_2}}$ yields
    $E_{A\cup B}=E_AE_B$; induct over the clusters partitioning $A$.
\end{proof}

\begin{lemma}[The operator of a cluster acts on its sites]
    \label{thm:hastings_cluster_operator_support}
    \lean{TNLean.Lattice.ClusterExpansion.clusterOp_mem_supportedOperators}
    \leanok
    \uses{def:hastings_cluster_operator, thm:hastings_trotter_product_form}
    If each $H_k$ acts on $S_k$, then $E_C$ acts on $\bigcup_{k\in C}S_k$.
\end{lemma}
\begin{proof}
    \leanok
    \uses{thm:hastings_trotter_product_form}
    For every $D\subseteq C$, $H_D$ and $e^{-\beta H_D}$ act on
    $\bigcup_{k\in C}S_k$.
\end{proof}

\begin{lemma}[The truncated operator in matrix product form]
    \label{thm:hastings_truncated_product_form}
    \lean{TNLean.Lattice.ClusterExpansion.IsLargeSupport,
        TNLean.Lattice.ClusterExpansion.noncommProd_apply_of_pairwiseDisjoint,
        TNLean.Lattice.ClusterExpansion.noLargeCluster_iff_filter,
        TNLean.Lattice.ClusterExpansion.truncatedExp_eq_mul,
        TNLean.Lattice.ClusterExpansion.sum_clusterOp_apply,
        TNLean.Lattice.ClusterExpansion.isLocalCoefficient_clusterCoeff,
        TNLean.Lattice.ClusterExpansion.sum_prod_clusterCoeff,
        TNLean.Lattice.ClusterExpansion.hasProductForm_sum_clusterOp,
        TNLean.Lattice.ClusterExpansion.hasProductForm_truncatedExp,
        TNLean.Lattice.ClusterExpansion.hasProductForm_truncatedExp_of_dist_le}
    \leanok
    \uses{def:hastings_product_form, def:hastings_truncated_expansion,
        def:hastings_cluster_of_terms}
    Let the domain be nonempty, let every term $H_k$ act only on a set $S_k$ of
    sites, let two terms be adjacent when their sets of sites meet, and call a
    set of terms large when its sites contain two sites at distance larger
    than $l$. Then the truncated operator $\rho(\beta,l)$ is in matrix product
    form of range $l$. This is the matrix product form at high temperature
    of~\cite[lines 1256--1279]{Hastings2006Solving}, on every finite metric
    lattice. The local index count is given below.
\end{lemma}
\begin{proof}
    \leanok
    \uses{thm:hastings_truncation_cluster_products,
        thm:hastings_cluster_operator_support,
        thm:hastings_product_form_supported_radius, thm:hastings_product_form_mul}
    By Lemmas~\ref{thm:hastings_truncation_cluster_products}
    and~\ref{thm:hastings_cluster_operator_support},
    $\rho(\beta,l)=\sum_A\prod_{C\in\mathrm{clusters}(A)}E_C$, where empty-support
    terms factor out as a scalar multiple of $1$ and each remaining cluster $C$
    acts on $S_C=\bigcup_{k\in C}S_k$ with $\operatorname{diam}S_C\le l$ and
    pairwise disjoint $S_C$. Indexing each site by its containing cluster (if
    any) and matrix-unit labels $(a_j,b_j)$ expresses $\sum_A\prod_CE_C$ in the
    form (\ref{eq:hastings_mpf}) of range $l$; conclude by
    Lemma~\ref{thm:hastings_product_form_mul}.
\end{proof}

\begin{lemma}[Local index ranges of the truncated operator]
    \label{thm:hastings_truncated_index_count}
    \lean{TNLean.Lattice.ProductForm.restrict,
        TNLean.Lattice.ProductForm.card_index_restrict,
        TNLean.Lattice.ClusterExpansion.clusterProductForm,
        TNLean.Lattice.ClusterExpansion.clusterIndexSet,
        TNLean.Lattice.ClusterExpansion.localClusterProductForm,
        TNLean.Lattice.ClusterExpansion.coeff_clusterProductForm_eq_zero_of_not_mem_clusterIndexSet,
        TNLean.Lattice.ClusterExpansion.card_clusterIndexSet_le,
        TNLean.Lattice.ClusterExpansion.IsValidCluster.subset_centres_ball,
        TNLean.Lattice.ClusterExpansion.card_clusterIndexSet_le_ball,
        TNLean.Lattice.ClusterExpansion.card_index_localClusterProductForm_le,
        TNLean.Lattice.ClusterExpansion.exists_productForm_truncatedExp_card_index,
        TNLean.Lattice.ClusterExpansion.exists_productForm_truncatedExp_pow_card_index,
        TNLean.Lattice.ClusterExpansion.exists_productForm_truncatedExp_card_index_of_polynomial_growth,
        TNLean.Lattice.ClusterExpansion.exists_productForm_truncatedExp_pow_card_index_of_polynomial_growth}
    \leanok
    \uses{thm:hastings_truncated_product_form, thm:hastings_product_form_mul}
    Let the site set be nonempty, let $q\ge1$, and suppose that each
    $H_k$ is supported within distance $R\ge0$ of its centre $c_k$. For $l\ge0$, put
    \begin{align*}
        M_i=\#\{k:d(i,c_k)\le l+R\}.
    \end{align*}
    The actual diameter-truncated operator $\rho(\beta,l)$ has a matrix
    product form of range $l$ with site index sets satisfying
    \begin{align*}
        |\mathcal A_i|\le q^2(1+2^{M_i}).
    \end{align*}
    Its $m$-th power has index ranges at most
    $[q^2(1+2^{M_i})]^m$; if the terms are Hermitian and $m$ is even,
    this power is positive semidefinite. If, moreover,
    $M_i\le C(\lceil l+R\rceil+1)^D$ for every $i$, these bounds become
    \begin{align*}
        q^2\bigl(1+2^{C(\lceil l+R\rceil+1)^D}\bigr),\qquad
        \bigl[q^2\bigl(1+2^{C(\lceil l+R\rceil+1)^D}\bigr)\bigr]^m.
    \end{align*}
    This makes explicit the polynomial volume growth used in the source's
    exponential dimensional estimates
    \cite[lines 1273--1275 and 1318--1320]{Hastings2006Solving}; a metric
    alone does not imply such growth\cite{gap:hastings05_planar_domain_scope}.
\end{lemma}
\begin{proof}\leanok
    Delete index values that record a cluster not containing their site;
    its coefficient is zero at that site. A retained cluster containing $i$
    has diameter at most $l$. Each of its terms acts on some site $j$, so
    $d(i,c_k)\le d(i,j)+d(j,c_k)\le l+R$. There are at most $2^{M_i}$
    possible sets of these terms. The additional value records no cluster,
    and each choice has $q^2$ matrix-unit labels. Empty-support terms split
    off as a scalar factor whose index sets have one element.
    Products multiply the index cardinalities, giving the bound for powers.
\end{proof}

\begin{lemma}[Polynomial ball growth in the integer lattice]
    \label{lem:hastings_integer_lattice_ball_count}
    \lean{TNLean.Lattice.card_centres_integerLattice_ball_le,
        TNLean.Lattice.card_centres_integerLattice_ball_le_polynomial,
        TNLean.Lattice.card_centres_ball_le_polynomial_of_dist_le,
        TNLean.Lattice.card_centres_ball_le_polynomial_of_isometry}
    \leanok
    For distinct centres in $\mathbb Z^D$ with the coordinate maximum
    metric, every radius-$r$ ball contains at most
    $(2\lceil r\rceil+1)^D\le2^D(\lceil r\rceil+1)^D$ centres.
    The same bound holds for every metric on a finite integer-lattice
    domain that dominates coordinate maximum distance, including its
    intrinsic walk metric. More generally, it transfers through any
    injective coordinate map that does not increase distances.
    Thus Lemma~\ref{thm:hastings_truncated_index_count} applies with
    $C=2^D$ on these lattices.
\end{lemma}
\begin{proof}\leanok
    Every coordinate of a centre in the ball lies in an integer interval
    with $2\lceil r\rceil+1$ elements. Containment in the product of these
    intervals proves the first bound. The second follows from
    $2\lceil r\rceil+1\le2(\lceil r\rceil+1)$.
\end{proof}
